\documentclass[11pt]{article}
\usepackage[margin=0.75in]{geometry}
\usepackage{fancyhdr}
\usepackage{array}
\usepackage{booktabs}
\usepackage{xcolor}
\usepackage{enumitem}

\pagestyle{fancy}
\fancyhf{}
\rhead{Patient: Santos, Maria E.}
\lhead{ProMotion Physical Therapy}
\rfoot{Page \thepage}

\definecolor{ptgreen}{RGB}{0,128,0}

\begin{document}

\begin{center}
{\Large\bfseries\color{ptgreen} PROMOTION PHYSICAL THERAPY}\\[0.3em]
{\small 900 Wellness Way, Suite 150, Riverside, CA 92501}\\
{\small Phone: (951) 555-7800 | Fax: (951) 555-7801}\\[1em]
{\LARGE\bfseries PHYSICAL THERAPY EVALUATION}
\end{center}

\vspace{0.5em}
\hrule height 2pt
\vspace{1em}

\section*{PATIENT INFORMATION}

\begin{tabular}{@{}p{3.5cm}p{5cm}p{3cm}p{4cm}@{}}
\textbf{Patient:} & Maria Elena Santos & \textbf{DOB:} & 03/14/1988 \\
\textbf{Eval Date:} & February 19, 2026 & \textbf{Diagnosis:} & Lumbar strain, \\
\textbf{Referring MD:} & Kevin Park, MD & & Post-concussion \\
\textbf{Therapist:} & Jennifer Martinez, DPT & \textbf{Insurance:} & Hartford WC \\
\end{tabular}

\section*{HISTORY}

35-year-old female carpenter referred for physical therapy following workplace injury on 02/08/2026. Patient fell approximately 15 feet from scaffolding, sustaining right distal radius fracture (currently in cast), concussion, and lumbar strain with MRI-confirmed L4-L5 disc herniation.

\textbf{Current Complaints:}
\begin{itemize}[nosep]
\item Low back pain rated 6/10, constant, aching quality
\item Pain radiates to left buttock and posterior thigh
\item Difficulty sitting for more than 15-20 minutes
\item Difficulty sleeping due to back pain
\item Stiffness in the morning
\item Residual headaches and dizziness (improving)
\end{itemize}

\section*{OBJECTIVE FINDINGS}

\textbf{Posture:} Guarded posture with slight forward lean. Decreased lumbar lordosis.

\textbf{Gait:} Antalgic gait pattern with shortened stride length bilaterally.

\textbf{Lumbar Range of Motion:}
\begin{tabular}{@{}lcc@{}}
\toprule
\textbf{Motion} & \textbf{Measured} & \textbf{Normal} \\
\midrule
Flexion & 40° & 80° \\
Extension & 10° & 30° \\
Left Lateral Flexion & 15° & 35° \\
Right Lateral Flexion & 18° & 35° \\
Left Rotation & 20° & 45° \\
Right Rotation & 22° & 45° \\
\bottomrule
\end{tabular}

\textbf{Palpation:} Tenderness and spasm of bilateral paraspinals L3-S1, more pronounced on left. Tenderness over L4-L5 spinous processes.

\textbf{Special Tests:}
\begin{itemize}[nosep]
\item Straight Leg Raise: Positive on left at 45° (reproduces radicular symptoms)
\item Slump Test: Positive on left
\item Prone Instability Test: Negative
\item FABER: Negative bilaterally
\end{itemize}

\textbf{Neurological:}
\begin{itemize}[nosep]
\item Sensation: Decreased to light touch left L5 dermatome (lateral leg, dorsum of foot)
\item Reflexes: 2+ bilateral patellar and Achilles
\item Motor: 4/5 left great toe extension (EHL weakness)
\end{itemize}

\textbf{Functional Assessment:}
\begin{itemize}[nosep]
\item Oswestry Disability Index: 48\% (Severe disability)
\item Unable to lift more than 5 lbs from floor
\item Unable to sit more than 20 minutes
\item Unable to stand more than 15 minutes
\end{itemize}

\section*{ASSESSMENT}

\begin{enumerate}
\item Lumbar disc herniation L4-L5 with left L5 radiculopathy
\item Lumbar paraspinal muscle strain/spasm
\item Decreased lumbar mobility and core stability
\item Functional limitations affecting ADLs and work capacity
\end{enumerate}

\section*{TREATMENT PLAN}

\textbf{Frequency:} 3x per week for 12 weeks (36 visits authorized)

\textbf{Goals (12 weeks):}
\begin{enumerate}
\item Reduce pain to 2/10 or less
\item Restore lumbar ROM to at least 75\% of normal
\item Improve Oswestry score to below 20\% (minimal disability)
\item Return to modified work duties
\item Independent with home exercise program
\end{enumerate}

\textbf{Treatment Interventions:}
\begin{itemize}[nosep]
\item Manual therapy: Soft tissue mobilization, joint mobilization
\item Therapeutic exercises: Core stabilization, McKenzie approach
\item Modalities: Heat, electrical stimulation for pain/spasm
\item Neuromuscular re-education
\item Patient education: Body mechanics, posture, activity modification
\item Progressive strengthening program
\end{itemize}

\vspace{1em}
\hrule
\vspace{1em}

\begin{center}
{\Large\bfseries PROGRESS NOTE - 6 WEEK UPDATE}\\
{\normalsize Date: April 1, 2026 | Visit \#18}
\end{center}

\vspace{1em}

\section*{SUBJECTIVE}

Patient reports steady improvement. Back pain decreased to 3-4/10. Radiating leg symptoms have significantly improved - now only occasional tingling with prolonged sitting. Sleeping better. Headaches have resolved. Right wrist cast was removed and she has begun occupational therapy for wrist as well.

\section*{OBJECTIVE}

\textbf{Lumbar ROM (Re-assessment):}
\begin{tabular}{@{}lccc@{}}
\toprule
\textbf{Motion} & \textbf{Initial} & \textbf{Current} & \textbf{Normal} \\
\midrule
Flexion & 40° & 65° & 80° \\
Extension & 10° & 22° & 30° \\
L Lateral Flex & 15° & 28° & 35° \\
\bottomrule
\end{tabular}

\textbf{Special Tests:}
\begin{itemize}[nosep]
\item SLR: Negative bilaterally (improved from positive at 45°)
\item Sensation: Improved - now intact L5 dermatome
\item Motor: 5/5 EHL (improved from 4/5)
\end{itemize}

\textbf{Functional:}
\begin{itemize}[nosep]
\item Oswestry: 28\% (improved from 48\%)
\item Can sit 45 minutes
\item Can lift 15 lbs safely
\end{itemize}

\section*{ASSESSMENT}

Patient making good progress. Radicular symptoms largely resolved. Core strength improving. On track to meet goals.

\section*{PLAN}

Continue PT 3x/week. Progress to more aggressive strengthening. Begin work conditioning activities. Coordinate with orthopedist regarding wrist recovery timeline.

\vspace{2em}

\begin{tabular}{@{}p{7cm}p{7cm}@{}}
\hrulefill & \\
\textbf{Jennifer Martinez, DPT} & \textbf{Date:} April 1, 2026 \\
Doctor of Physical Therapy & \\
ProMotion Physical Therapy & \\
\end{tabular}

\end{document}
